\subsection{Quotient by a closed two-sided ideal}

Lemma 15. --- Let A be a stellar algebra and let I be a closed two-sided ideal of A. Then I is self-adjoint.

Let \(J = I \cap I^*\) . The set \(J\) is a self-adjoint two-sided ideal of \(A\) , which contains \(I^*I\) . In particular, \(J\) is a stellar algebra. Let \(\mathfrak{F}\) be an increasing approximate unit of \(J\) (prop. 18 of I, p.~121). For all \(x \in I\) and every \(f \in J_{+}^{\leqslant 1}\) , one has

\[
\begin{array}{c} \| x f - x \| ^ {2} = \| (x f - x) ^ {*} (x f - x) \| = \| f (x ^ {*} x f - x ^ {*} x) - (x ^ {*} x f - x ^ {*} x) \| \\ \leqslant 2 \| x ^ {*} x f - x ^ {*} x \|. \end{array}
\]

Since \(x^{*}x\in \mathrm{J}\) , we therefore have

\[
\lim _ {f, \mathfrak {F}} \| x f - x \| ^ {2} = 0.
\]

Since \(xf \in \mathrm{J}\) for all \(f \in \mathrm{J}_{+}^{\leqslant 1}\) and since \(\mathrm{J}\) is closed, this implies that \(x \in \mathrm{J}\) . Therefore \(\mathrm{I} = \mathrm{J}\) , and the ideal \(\mathrm{I}\) is self-adjoint.

PROPOSITION 19. --- Let A be a stellar algebra and let I be a closed two-sided ideal of A. Then the quotient involutive Banach algebra A/I is a stellar algebra.

The ideal I is self-adjoint (Lemma 15). We consider the stellar algebra \(\widetilde{\mathsf{A}}\) deduced from A by adjoining an identity element (def. 4 of I, p.~106). In this algebra, the set I is a closed self-adjoint two-sided ideal, and A/I is identified with a closed involutive subalgebra of \(\widetilde{\mathrm{A}}/\mathrm{I}\) . We may therefore assume that A is unital.

The Banach algebra A/I is involutive. Let \(\pi\colon\mathrm{A}\to\mathrm{A}/\mathrm{I}\) be the canonical projection. The self-adjoint two-sided ideal I is a stellar subalgebra of A. Let \(\mathfrak{F}\) be an increasing approximate unit of I (prop. 18 of I, p.~121). Let us first show that for every \(x\in\mathrm{A}\) , one has

\[
\| \pi (x) \| _ {\mathrm{A/I}} = \lim _ {f, \mathfrak {F}} \| x - x f \|.\tag{6}
\]

On the one hand, since \(xf \in \mathrm{I}\) for all \(f \in \mathrm{I}_{+}^{\leqslant 1}\) , one has

\[
\| \pi (x) \| _ {\mathrm{A} / \mathrm{I}} = \inf _ {a \in \mathrm{I}} \| x - a \| \leqslant \operatorname * {l i m i n f} _ {f, \mathfrak {F}} \| x - x f \|.
\]

On the other hand, for every \(a \in \mathrm{I}\) , one has

\[
\begin{array}{c} \| x - x f \| \leqslant \| (x - a) - (x - a) f \| + \| a - a f \| \\ = \| (x - a) (1 - f) \| + \| a - a f \| \end{array}
\]

and therefore, since \(\|1 - f\| \leqslant 1\) (Lemma 12 of I, p.~116) and \(a \in I\), from which we deduce

\[
\operatorname * {l i m s u p} _ {f, \mathfrak {F}} \| x - x f \| \leqslant \| x - a \|.
\]

Thus, it follows \(\limsup_{f,\mathfrak{F}} \|x - xf\| \leqslant \|\pi(x)\|_{\mathrm{A/I}}\) since \(a\) is arbitrary in I. Formula (6) is therefore proved.

Now let \(x\) be an element of A. By formula (6), we have

\[
\begin{array}{c} \| \pi (x) \| _ {\mathrm{A/I}} ^ {2} = \lim _ {f, \mathfrak {F}} \| x - x f \| ^ {2} = \lim _ {f, \mathfrak {F}} \| x (1 - f) \| ^ {2} \\ = \lim _ {f, \mathfrak {F}} \| (1 - f) x ^ {*} x (1 - f) \| \leqslant \lim _ {f, \mathfrak {F}} \| x ^ {*} x (1 - f) \| = \| \pi (x ^ {*} x) \| _ {\mathrm{A/I}}. \end{array}
\]

Lemma 6 of I, p.~103 then implies that the involutive Banach algebra A/I is a stellar algebra.

